%% Generated by Sphinx.
\def\sphinxdocclass{report}
\documentclass[letterpaper,10pt,english]{sphinxmanual}
\ifdefined\pdfpxdimen
   \let\sphinxpxdimen\pdfpxdimen\else\newdimen\sphinxpxdimen
\fi \sphinxpxdimen=.75bp\relax
\ifdefined\pdfimageresolution
    \pdfimageresolution= \numexpr \dimexpr1in\relax/\sphinxpxdimen\relax
\fi
\newdimen\sphinxremdimen\sphinxremdimen = 10pt
%% let collapsible pdf bookmarks panel have high depth per default
\PassOptionsToPackage{bookmarksdepth=5}{hyperref}

\PassOptionsToPackage{booktabs}{sphinx}
\PassOptionsToPackage{colorrows}{sphinx}

\PassOptionsToPackage{warn}{textcomp}
\usepackage[utf8]{inputenc}
\ifdefined\DeclareUnicodeCharacter
% support both utf8 and utf8x syntaxes
  \ifdefined\DeclareUnicodeCharacterAsOptional
    \def\sphinxDUC#1{\DeclareUnicodeCharacter{"#1}}
  \else
    \let\sphinxDUC\DeclareUnicodeCharacter
  \fi
  \sphinxDUC{00A0}{\nobreakspace}
  \sphinxDUC{2500}{\sphinxunichar{2500}}
  \sphinxDUC{2502}{\sphinxunichar{2502}}
  \sphinxDUC{2514}{\sphinxunichar{2514}}
  \sphinxDUC{251C}{\sphinxunichar{251C}}
  \sphinxDUC{2572}{\textbackslash}
\fi
\usepackage{cmap}
\usepackage[T1]{fontenc}
\usepackage{amsmath,amssymb,amstext}
\usepackage{babel}



\usepackage{tgtermes}
\usepackage{tgheros}
\renewcommand{\ttdefault}{txtt}



\usepackage[Bjarne]{fncychap}
\usepackage{sphinx}

\fvset{fontsize=auto}
\usepackage{geometry}


% Include hyperref last.
\usepackage{hyperref}
% Fix anchor placement for figures with captions.
\usepackage{hypcap}% it must be loaded after hyperref.
% Set up styles of URL: it should be placed after hyperref.
\urlstyle{same}

\addto\captionsenglish{\renewcommand{\contentsname}{Contents:}}

\usepackage{sphinxmessages}
\setcounter{tocdepth}{0}



\title{dtw\_loss\_functions}
\date{Mar 06, 2026}
\release{}
\author{Alberto (Jesus) Zancanaro}
\newcommand{\sphinxlogo}{\vbox{}}
\renewcommand{\releasename}{}
\makeindex
\begin{document}

\ifdefined\shorthandoff
  \ifnum\catcode`\=\string=\active\shorthandoff{=}\fi
  \ifnum\catcode`\"=\active\shorthandoff{"}\fi
\fi

\pagestyle{empty}
\sphinxmaketitle
\pagestyle{plain}
\sphinxtableofcontents
\pagestyle{normal}
\phantomsection\label{\detokenize{index::doc}}


\sphinxAtStartPar
This documentation is still work in progress. If you have any questions, please feel free to contact the author.
For now the content is generated automatically from the docstrings in the code, through sphinx autodoc. The documentation will be improved in the future, with more examples and explanations.

\sphinxAtStartPar
\sphinxcode{\sphinxupquote{dtw\_loss\_functions}} is an python package with the torch implementation of various DTW\sphinxhyphen{}derived loss functions.

\sphinxstepscope


\chapter{Block DTW}
\label{\detokenize{dtw_loss_functions.block_dtw:module-dtw_loss_functions.block_dtw}}\label{\detokenize{dtw_loss_functions.block_dtw:block-dtw}}\label{\detokenize{dtw_loss_functions.block_dtw::doc}}\index{module@\spxentry{module}!dtw\_loss\_functions.block\_dtw@\spxentry{dtw\_loss\_functions.block\_dtw}}\index{dtw\_loss\_functions.block\_dtw@\spxentry{dtw\_loss\_functions.block\_dtw}!module@\spxentry{module}}
\sphinxAtStartPar
Implementation of the block DTW, which is a variant of the SDTW that computes the SDTW on blocks of the signal instead of the entire signal.
\subsubsection*{Example}

\begin{sphinxVerbatim}[commandchars=\\\{\}]
\PYG{g+gp}{\PYGZgt{}\PYGZgt{}\PYGZgt{} }\PYG{k+kn}{from}\PYG{+w}{ }\PYG{n+nn}{dtw\PYGZus{}loss\PYGZus{}functions}\PYG{+w}{ }\PYG{k+kn}{import} \PYG{n}{block\PYGZus{}dtw}
\PYG{g+gp}{\PYGZgt{}\PYGZgt{}\PYGZgt{} }\PYG{k+kn}{import}\PYG{+w}{ }\PYG{n+nn}{torch}
\PYG{g+gp}{\PYGZgt{}\PYGZgt{}\PYGZgt{} }\PYG{n}{block\PYGZus{}size} \PYG{o}{=} \PYG{l+m+mi}{25}
\PYG{g+gp}{\PYGZgt{}\PYGZgt{}\PYGZgt{} }\PYG{n}{use\PYGZus{}cuda} \PYG{o}{=} \PYG{n}{torch}\PYG{o}{.}\PYG{n}{cuda}\PYG{o}{.}\PYG{n}{is\PYGZus{}available}\PYG{p}{(}\PYG{p}{)}
\PYG{g+gp}{\PYGZgt{}\PYGZgt{}\PYGZgt{} }\PYG{n}{device} \PYG{o}{=} \PYG{l+s+s1}{\PYGZsq{}}\PYG{l+s+s1}{cuda}\PYG{l+s+s1}{\PYGZsq{}} \PYG{k}{if} \PYG{n}{use\PYGZus{}cuda} \PYG{k}{else} \PYG{l+s+s1}{\PYGZsq{}}\PYG{l+s+s1}{cpu}\PYG{l+s+s1}{\PYGZsq{}}
\PYG{g+gp}{\PYGZgt{}\PYGZgt{}\PYGZgt{} }\PYG{n}{batch\PYGZus{}size} \PYG{o}{=} \PYG{l+m+mi}{5}
\PYG{g+gp}{\PYGZgt{}\PYGZgt{}\PYGZgt{} }\PYG{n}{time\PYGZus{}samples} \PYG{o}{=} \PYG{l+m+mi}{300}
\PYG{g+gp}{\PYGZgt{}\PYGZgt{}\PYGZgt{} }\PYG{n}{channels} \PYG{o}{=} \PYG{l+m+mi}{1}
\PYG{g+gp}{\PYGZgt{}\PYGZgt{}\PYGZgt{} }\PYG{n}{x}   \PYG{o}{=} \PYG{n}{torch}\PYG{o}{.}\PYG{n}{randn}\PYG{p}{(}\PYG{n}{batch\PYGZus{}size}\PYG{p}{,} \PYG{n}{time\PYGZus{}samples}\PYG{p}{,} \PYG{n}{channels}\PYG{p}{)}\PYG{o}{.}\PYG{n}{to}\PYG{p}{(}\PYG{n}{device}\PYG{p}{)}
\PYG{g+gp}{\PYGZgt{}\PYGZgt{}\PYGZgt{} }\PYG{n}{x\PYGZus{}r} \PYG{o}{=} \PYG{n}{torch}\PYG{o}{.}\PYG{n}{randn}\PYG{p}{(}\PYG{n}{batch\PYGZus{}size}\PYG{p}{,} \PYG{n}{time\PYGZus{}samples}\PYG{p}{,} \PYG{n}{channels}\PYG{p}{)}\PYG{o}{.}\PYG{n}{to}\PYG{p}{(}\PYG{n}{device}\PYG{p}{)}
\PYG{g+gp}{\PYGZgt{}\PYGZgt{}\PYGZgt{} }\PYG{n}{block\PYGZus{}dtw\PYGZus{}loss} \PYG{o}{=} \PYG{n}{block\PYGZus{}dtw}\PYG{o}{.}\PYG{n}{block\PYGZus{}dtw}\PYG{p}{(}\PYG{n}{block\PYGZus{}size}\PYG{p}{,} \PYG{n}{use\PYGZus{}cuda}\PYG{p}{)}
\PYG{g+gp}{\PYGZgt{}\PYGZgt{}\PYGZgt{} }\PYG{n}{output\PYGZus{}block\PYGZus{}dtw} \PYG{o}{=} \PYG{n}{block\PYGZus{}dtw\PYGZus{}loss}\PYG{p}{(}\PYG{n}{x}\PYG{p}{,} \PYG{n}{x\PYGZus{}r}\PYG{p}{)}
\end{sphinxVerbatim}


\section{Authors}
\label{\detokenize{dtw_loss_functions.block_dtw:authors}}
\sphinxAtStartPar
Alberto Zancanaro \textless{}\sphinxhref{mailto:alberto.zancanaro@uni.lu}{alberto.zancanaro@uni.lu}\textgreater{}
\index{block\_dtw (class in dtw\_loss\_functions.block\_dtw)@\spxentry{block\_dtw}\spxextra{class in dtw\_loss\_functions.block\_dtw}}

\begin{fulllineitems}
\phantomsection\label{\detokenize{dtw_loss_functions.block_dtw:dtw_loss_functions.block_dtw.block_dtw}}
\pysigstartsignatures
\pysiglinewithargsret
{\sphinxbfcode{\sphinxupquote{\DUrole{k}{class}\DUrole{w}{ }}}\sphinxcode{\sphinxupquote{dtw\_loss\_functions.block\_dtw.}}\sphinxbfcode{\sphinxupquote{block\_dtw}}}
{\sphinxparam{\DUrole{n}{block\_size}\DUrole{p}{:}\DUrole{w}{ }\DUrole{n}{int}}\sphinxparamcomma \sphinxparam{\DUrole{n}{use\_cuda}\DUrole{p}{:}\DUrole{w}{ }\DUrole{n}{bool}}\sphinxparamcomma \sphinxparam{\DUrole{n}{gamma\_sdtw}\DUrole{p}{:}\DUrole{w}{ }\DUrole{n}{float}\DUrole{w}{ }\DUrole{o}{=}\DUrole{w}{ }\DUrole{default_value}{1}}\sphinxparamcomma \sphinxparam{\DUrole{n}{use\_divergence}\DUrole{p}{:}\DUrole{w}{ }\DUrole{n}{bool}\DUrole{w}{ }\DUrole{o}{=}\DUrole{w}{ }\DUrole{default_value}{False}}\sphinxparamcomma \sphinxparam{\DUrole{n}{bandwidth}\DUrole{p}{:}\DUrole{w}{ }\DUrole{n}{float}\DUrole{w}{ }\DUrole{o}{=}\DUrole{w}{ }\DUrole{default_value}{None}}\sphinxparamcomma \sphinxparam{\DUrole{n}{dist\_func}\DUrole{o}{=}\DUrole{default_value}{None}}}
{}
\pysigstopsignatures
\sphinxAtStartPar
Bases: \sphinxcode{\sphinxupquote{object}}

\sphinxAtStartPar
Class that compute the block DTW loss, which is a variant of the SDTW that computes the SDTW on blocks of the signal instead of the entire signal.
The block DTW can be computed in two ways:
\begin{itemize}
\item {} 
\sphinxAtStartPar
Naive (SEQUENTIAL) implementation: compute the SDTW on each block separately and sum the results.

\item {} 
\sphinxAtStartPar
Optimized (PARALLEL) implementation: exploit reshaping of the input tensors to compute the SDTW on all blocks at once.

\end{itemize}

\sphinxAtStartPar
This class will select automatically which implementation to use based on the input tensors length and block size.
See the docstring of block\_dtw\_optimized for more details on the requirements for the optimized implementation.

\sphinxAtStartPar
If you are not sure which implementation to use, you can use the block\_dtw class.
Note that if you know a priori that the optimized version can be used in your case, it is recommended to use directly the block\_dtw\_optimized class, which is faster than the block\_dtw class (no overhead of checking the input tensors length and block size).
\index{block\_size (dtw\_loss\_functions.block\_dtw.block\_dtw attribute)@\spxentry{block\_size}\spxextra{dtw\_loss\_functions.block\_dtw.block\_dtw attribute}}

\begin{fulllineitems}
\phantomsection\label{\detokenize{dtw_loss_functions.block_dtw:dtw_loss_functions.block_dtw.block_dtw.block_size}}
\pysigstartsignatures
\pysigline
{\sphinxbfcode{\sphinxupquote{block\_size}}}
\pysigstopsignatures
\sphinxAtStartPar
Size of the blocks into which to divide the signal.
\begin{quote}\begin{description}
\sphinxlineitem{Type}
\sphinxAtStartPar
int

\end{description}\end{quote}

\end{fulllineitems}

\index{block\_dtw\_naive (dtw\_loss\_functions.block\_dtw.block\_dtw attribute)@\spxentry{block\_dtw\_naive}\spxextra{dtw\_loss\_functions.block\_dtw.block\_dtw attribute}}

\begin{fulllineitems}
\phantomsection\label{\detokenize{dtw_loss_functions.block_dtw:dtw_loss_functions.block_dtw.block_dtw.block_dtw_naive}}
\pysigstartsignatures
\pysigline
{\sphinxbfcode{\sphinxupquote{block\_dtw\_naive}}}
\pysigstopsignatures
\sphinxAtStartPar
Instance of the naive implementation of the block DTW.
\begin{quote}\begin{description}
\sphinxlineitem{Type}
\sphinxAtStartPar
\sphinxcode{\sphinxupquote{block\_dtw\_naive}}

\end{description}\end{quote}

\end{fulllineitems}

\index{block\_dtw\_optimized (dtw\_loss\_functions.block\_dtw.block\_dtw attribute)@\spxentry{block\_dtw\_optimized}\spxextra{dtw\_loss\_functions.block\_dtw.block\_dtw attribute}}

\begin{fulllineitems}
\phantomsection\label{\detokenize{dtw_loss_functions.block_dtw:dtw_loss_functions.block_dtw.block_dtw.block_dtw_optimized}}
\pysigstartsignatures
\pysigline
{\sphinxbfcode{\sphinxupquote{block\_dtw\_optimized}}}
\pysigstopsignatures
\sphinxAtStartPar
Instance of the optimized implementation of the block DTW.
\begin{quote}\begin{description}
\sphinxlineitem{Type}
\sphinxAtStartPar
\sphinxcode{\sphinxupquote{block\_dtw\_optimized}}

\end{description}\end{quote}

\end{fulllineitems}

\begin{quote}\begin{description}
\sphinxlineitem{Parameters}\begin{itemize}
\item {} 
\sphinxAtStartPar
\sphinxstyleliteralstrong{\sphinxupquote{block\_size}} (\sphinxstyleliteralemphasis{\sphinxupquote{int}}) \textendash{} Size of the blocks into which to divide the signal.

\item {} 
\sphinxAtStartPar
\sphinxstyleliteralstrong{\sphinxupquote{use\_cuda}} (\sphinxstyleliteralemphasis{\sphinxupquote{bool}}) \textendash{} If true, this class will use the CUDA implementation of the SDTW.

\item {} 
\sphinxAtStartPar
\sphinxstyleliteralstrong{\sphinxupquote{gamma\_sdtw}} (\sphinxstyleliteralemphasis{\sphinxupquote{float}}\sphinxstyleliteralemphasis{\sphinxupquote{, }}\sphinxstyleliteralemphasis{\sphinxupquote{optional}}) \textendash{} Value of the gamma hyperparameter for the SDTW. Default is 1.

\item {} 
\sphinxAtStartPar
\sphinxstyleliteralstrong{\sphinxupquote{use\_divergence}} (\sphinxstyleliteralemphasis{\sphinxupquote{bool}}\sphinxstyleliteralemphasis{\sphinxupquote{, }}\sphinxstyleliteralemphasis{\sphinxupquote{optional}}) \textendash{} If true, compute the SDTW divergence instead of the SDTW. Default is False.

\item {} 
\sphinxAtStartPar
\sphinxstyleliteralstrong{\sphinxupquote{bandwidth}} (\sphinxstyleliteralemphasis{\sphinxupquote{float}}\sphinxstyleliteralemphasis{\sphinxupquote{, }}\sphinxstyleliteralemphasis{\sphinxupquote{optional}}) \textendash{} Sakoe\sphinxhyphen{}Chiba bandwidth for pruning. If the ‘None’ is given, no pruning is applied. Default is None.

\item {} 
\sphinxAtStartPar
\sphinxstyleliteralstrong{\sphinxupquote{dist\_func}} (\sphinxstyleliteralemphasis{\sphinxupquote{function}}\sphinxstyleliteralemphasis{\sphinxupquote{, }}\sphinxstyleliteralemphasis{\sphinxupquote{optional}}) \textendash{} Distance function to use for the SDTW. Default is None, which corresponds to the squared Euclidean distance.

\end{itemize}

\end{description}\end{quote}
\subsubsection*{Methods}


\begin{savenotes}
\sphinxatlongtablestart
\sphinxthistablewithglobalstyle
\sphinxthistablewithnovlinesstyle
\makeatletter
  \LTleft \@totalleftmargin plus1fill
  \LTright\dimexpr\columnwidth-\@totalleftmargin-\linewidth\relax plus1fill
\makeatother
\begin{longtable}{\X{1}{2}\X{1}{2}}
\sphinxtoprule
\endfirsthead

\multicolumn{2}{c}{\sphinxnorowcolor
    \makebox[0pt]{\sphinxtablecontinued{\tablename\ \thetable{} \textendash{} continued from previous page}}%
}\\
\sphinxtoprule
\endhead

\sphinxbottomrule
\multicolumn{2}{r}{\sphinxnorowcolor
    \makebox[0pt][r]{\sphinxtablecontinued{continues on next page}}%
}\\
\endfoot

\endlastfoot
\sphinxtableatstartofbodyhook
\begin{varwidth}[t]{\sphinxcolwidth{1}{2}}
\sphinxAtStartPar
\sphinxcode{\sphinxupquote{\_\_call\_\_}}(x, x\_r)
\sphinxbeforeendvarwidth
\end{varwidth}%
&\begin{varwidth}[t]{\sphinxcolwidth{1}{2}}
\sphinxAtStartPar
Compute the block DTW loss between the input tensors \sphinxtitleref{x} and \sphinxtitleref{x\_r}.
\sphinxbeforeendvarwidth
\end{varwidth}%
\\
\sphinxbottomrule
\end{longtable}
\sphinxtableafterendhook
\sphinxatlongtableend
\end{savenotes}

\end{fulllineitems}

\index{block\_dtw\_naive (class in dtw\_loss\_functions.block\_dtw)@\spxentry{block\_dtw\_naive}\spxextra{class in dtw\_loss\_functions.block\_dtw}}

\begin{fulllineitems}
\phantomsection\label{\detokenize{dtw_loss_functions.block_dtw:dtw_loss_functions.block_dtw.block_dtw_naive}}
\pysigstartsignatures
\pysiglinewithargsret
{\sphinxbfcode{\sphinxupquote{\DUrole{k}{class}\DUrole{w}{ }}}\sphinxcode{\sphinxupquote{dtw\_loss\_functions.block\_dtw.}}\sphinxbfcode{\sphinxupquote{block\_dtw\_naive}}}
{\sphinxparam{\DUrole{n}{block\_size}\DUrole{p}{:}\DUrole{w}{ }\DUrole{n}{int}}\sphinxparamcomma \sphinxparam{\DUrole{n}{use\_cuda}\DUrole{p}{:}\DUrole{w}{ }\DUrole{n}{bool}}\sphinxparamcomma \sphinxparam{\DUrole{n}{gamma\_sdtw}\DUrole{p}{:}\DUrole{w}{ }\DUrole{n}{float}\DUrole{w}{ }\DUrole{o}{=}\DUrole{w}{ }\DUrole{default_value}{1}}\sphinxparamcomma \sphinxparam{\DUrole{n}{use\_divergence}\DUrole{p}{:}\DUrole{w}{ }\DUrole{n}{bool}\DUrole{w}{ }\DUrole{o}{=}\DUrole{w}{ }\DUrole{default_value}{False}}\sphinxparamcomma \sphinxparam{\DUrole{n}{bandwidth}\DUrole{p}{:}\DUrole{w}{ }\DUrole{n}{float}\DUrole{w}{ }\DUrole{o}{=}\DUrole{w}{ }\DUrole{default_value}{None}}\sphinxparamcomma \sphinxparam{\DUrole{n}{dist\_func}\DUrole{o}{=}\DUrole{default_value}{None}}}
{}
\pysigstopsignatures
\sphinxAtStartPar
Bases: \sphinxcode{\sphinxupquote{SoftDTW}}

\sphinxAtStartPar
Naive implementation of the block DTW, which computes the SDTW on each block separately.

\sphinxAtStartPar
For details on the parameters, see the docstring of the \sphinxcode{\sphinxupquote{block\_dtw}} class.
\subsubsection*{Methods}


\begin{savenotes}
\sphinxatlongtablestart
\sphinxthistablewithglobalstyle
\sphinxthistablewithnovlinesstyle
\makeatletter
  \LTleft \@totalleftmargin plus1fill
  \LTright\dimexpr\columnwidth-\@totalleftmargin-\linewidth\relax plus1fill
\makeatother
\begin{longtable}{\X{1}{2}\X{1}{2}}
\sphinxtoprule
\endfirsthead

\multicolumn{2}{c}{\sphinxnorowcolor
    \makebox[0pt]{\sphinxtablecontinued{\tablename\ \thetable{} \textendash{} continued from previous page}}%
}\\
\sphinxtoprule
\endhead

\sphinxbottomrule
\multicolumn{2}{r}{\sphinxnorowcolor
    \makebox[0pt][r]{\sphinxtablecontinued{continues on next page}}%
}\\
\endfoot

\endlastfoot
\sphinxtableatstartofbodyhook
\begin{varwidth}[t]{\sphinxcolwidth{1}{2}}
\sphinxAtStartPar
\sphinxcode{\sphinxupquote{forward}}(x, x\_r)
\sphinxbeforeendvarwidth
\end{varwidth}%
&\begin{varwidth}[t]{\sphinxcolwidth{1}{2}}
\sphinxAtStartPar
Compute the block DTW loss between the input tensors \sphinxtitleref{x} and \sphinxtitleref{x\_r} by computing the SDTW on each block separately and summing the results.
\sphinxbeforeendvarwidth
\end{varwidth}%
\\
\sphinxbottomrule
\end{longtable}
\sphinxtableafterendhook
\sphinxatlongtableend
\end{savenotes}
\index{forward() (dtw\_loss\_functions.block\_dtw.block\_dtw\_naive method)@\spxentry{forward()}\spxextra{dtw\_loss\_functions.block\_dtw.block\_dtw\_naive method}}

\begin{fulllineitems}
\phantomsection\label{\detokenize{dtw_loss_functions.block_dtw:dtw_loss_functions.block_dtw.block_dtw_naive.forward}}
\pysigstartsignatures
\pysiglinewithargsret
{\sphinxbfcode{\sphinxupquote{forward}}}
{\sphinxparam{\DUrole{n}{x}\DUrole{p}{:}\DUrole{w}{ }\DUrole{n}{tensor}}\sphinxparamcomma \sphinxparam{\DUrole{n}{x\_r}\DUrole{p}{:}\DUrole{w}{ }\DUrole{n}{tensor}}}
{{ $\rightarrow$ tensor}}
\pysigstopsignatures
\sphinxAtStartPar
Compute the block DTW loss between the input tensors \sphinxtitleref{x} and \sphinxtitleref{x\_r} by computing the SDTW on each block separately and summing the results.
\begin{quote}\begin{description}
\sphinxlineitem{Parameters}\begin{itemize}
\item {} 
\sphinxAtStartPar
\sphinxstyleliteralstrong{\sphinxupquote{x}} (\sphinxstyleliteralemphasis{\sphinxupquote{torch.tensor}}) \textendash{} First input tensor of shape B x T x C

\item {} 
\sphinxAtStartPar
\sphinxstyleliteralstrong{\sphinxupquote{x\_r}} (\sphinxstyleliteralemphasis{\sphinxupquote{torch.tensor}}) \textendash{} Second input tensor of shape B x T x C

\end{itemize}

\sphinxlineitem{Returns}
\sphinxAtStartPar
\sphinxstylestrong{recon\_error} \textendash{} Tensor of shape B containing the block DTW loss for each sample in the batch.

\sphinxlineitem{Return type}
\sphinxAtStartPar
torch.tensor

\end{description}\end{quote}

\end{fulllineitems}


\end{fulllineitems}

\index{block\_dtw\_optimized (class in dtw\_loss\_functions.block\_dtw)@\spxentry{block\_dtw\_optimized}\spxextra{class in dtw\_loss\_functions.block\_dtw}}

\begin{fulllineitems}
\phantomsection\label{\detokenize{dtw_loss_functions.block_dtw:dtw_loss_functions.block_dtw.block_dtw_optimized}}
\pysigstartsignatures
\pysiglinewithargsret
{\sphinxbfcode{\sphinxupquote{\DUrole{k}{class}\DUrole{w}{ }}}\sphinxcode{\sphinxupquote{dtw\_loss\_functions.block\_dtw.}}\sphinxbfcode{\sphinxupquote{block\_dtw\_optimized}}}
{\sphinxparam{\DUrole{n}{block\_size}\DUrole{p}{:}\DUrole{w}{ }\DUrole{n}{int}}\sphinxparamcomma \sphinxparam{\DUrole{n}{use\_cuda}\DUrole{p}{:}\DUrole{w}{ }\DUrole{n}{bool}}\sphinxparamcomma \sphinxparam{\DUrole{n}{gamma\_sdtw}\DUrole{p}{:}\DUrole{w}{ }\DUrole{n}{float}\DUrole{w}{ }\DUrole{o}{=}\DUrole{w}{ }\DUrole{default_value}{1}}\sphinxparamcomma \sphinxparam{\DUrole{n}{use\_divergence}\DUrole{p}{:}\DUrole{w}{ }\DUrole{n}{bool}\DUrole{w}{ }\DUrole{o}{=}\DUrole{w}{ }\DUrole{default_value}{False}}\sphinxparamcomma \sphinxparam{\DUrole{n}{bandwidth}\DUrole{p}{:}\DUrole{w}{ }\DUrole{n}{float}\DUrole{w}{ }\DUrole{o}{=}\DUrole{w}{ }\DUrole{default_value}{None}}\sphinxparamcomma \sphinxparam{\DUrole{n}{dist\_func}\DUrole{o}{=}\DUrole{default_value}{None}}}
{}
\pysigstopsignatures
\sphinxAtStartPar
Bases: \sphinxcode{\sphinxupquote{SoftDTW}}

\sphinxAtStartPar
Optimized implementation of the block DTW, which exploits reshaping of the input tensors to compute the SDTW on all blocks at once.

\sphinxAtStartPar
This version can be used only if the length of the input tensors is divisible by the block size, i.e. if \sphinxcode{\sphinxupquote{length\_signal \% block\_size == 0}}.
Note that the class will not check if this condition is satisfied, so it is the responsibility of the user to ensure that the input tensors length and block size are compatible.

\sphinxAtStartPar
This requirement is necessary because the optimized implementation exploits reshaping of the input tensors to compute the SDTW on all blocks at once.
This works because SoftDTW implementation allow batched inputs,so we can reshape the input tensors to have a new batch size equal to the number of blocks, and compute the SDTW on all blocks at once.

\sphinxAtStartPar
For details on the parameters, see the docstring of the \sphinxcode{\sphinxupquote{block\_dtw}} class.
\subsubsection*{Methods}


\begin{savenotes}
\sphinxatlongtablestart
\sphinxthistablewithglobalstyle
\sphinxthistablewithnovlinesstyle
\makeatletter
  \LTleft \@totalleftmargin plus1fill
  \LTright\dimexpr\columnwidth-\@totalleftmargin-\linewidth\relax plus1fill
\makeatother
\begin{longtable}{\X{1}{2}\X{1}{2}}
\sphinxtoprule
\endfirsthead

\multicolumn{2}{c}{\sphinxnorowcolor
    \makebox[0pt]{\sphinxtablecontinued{\tablename\ \thetable{} \textendash{} continued from previous page}}%
}\\
\sphinxtoprule
\endhead

\sphinxbottomrule
\multicolumn{2}{r}{\sphinxnorowcolor
    \makebox[0pt][r]{\sphinxtablecontinued{continues on next page}}%
}\\
\endfoot

\endlastfoot
\sphinxtableatstartofbodyhook
\begin{varwidth}[t]{\sphinxcolwidth{1}{2}}
\sphinxAtStartPar
\sphinxcode{\sphinxupquote{forward}}(x, x\_r)
\sphinxbeforeendvarwidth
\end{varwidth}%
&\begin{varwidth}[t]{\sphinxcolwidth{1}{2}}
\sphinxAtStartPar
Compute the block DTW loss between the input tensors \sphinxtitleref{x} and \sphinxtitleref{x\_r} by exploiting reshaping of the input tensors to compute the SDTW on all blocks at once.
\sphinxbeforeendvarwidth
\end{varwidth}%
\\
\sphinxbottomrule
\end{longtable}
\sphinxtableafterendhook
\sphinxatlongtableend
\end{savenotes}
\index{forward() (dtw\_loss\_functions.block\_dtw.block\_dtw\_optimized method)@\spxentry{forward()}\spxextra{dtw\_loss\_functions.block\_dtw.block\_dtw\_optimized method}}

\begin{fulllineitems}
\phantomsection\label{\detokenize{dtw_loss_functions.block_dtw:dtw_loss_functions.block_dtw.block_dtw_optimized.forward}}
\pysigstartsignatures
\pysiglinewithargsret
{\sphinxbfcode{\sphinxupquote{forward}}}
{\sphinxparam{\DUrole{n}{x}\DUrole{p}{:}\DUrole{w}{ }\DUrole{n}{tensor}}\sphinxparamcomma \sphinxparam{\DUrole{n}{x\_r}\DUrole{p}{:}\DUrole{w}{ }\DUrole{n}{tensor}}}
{{ $\rightarrow$ tensor}}
\pysigstopsignatures
\sphinxAtStartPar
Compute the block DTW loss between the input tensors \sphinxtitleref{x} and \sphinxtitleref{x\_r} by exploiting reshaping of the input tensors to compute the SDTW on all blocks at once.
\begin{quote}\begin{description}
\sphinxlineitem{Parameters}\begin{itemize}
\item {} 
\sphinxAtStartPar
\sphinxstyleliteralstrong{\sphinxupquote{x}} (\sphinxstyleliteralemphasis{\sphinxupquote{torch.tensor}}) \textendash{} First input tensor of shape B x T x C

\item {} 
\sphinxAtStartPar
\sphinxstyleliteralstrong{\sphinxupquote{x\_r}} (\sphinxstyleliteralemphasis{\sphinxupquote{torch.tensor}}) \textendash{} Second input tensor of shape B x T x C

\end{itemize}

\sphinxlineitem{Returns}
\sphinxAtStartPar
\sphinxstylestrong{recon\_error} \textendash{} Tensor of shape B containing the block DTW loss for each sample in the batch.

\sphinxlineitem{Return type}
\sphinxAtStartPar
torch.tensor

\end{description}\end{quote}

\end{fulllineitems}


\end{fulllineitems}


\sphinxstepscope


\chapter{OTW}
\label{\detokenize{dtw_loss_functions.otw:module-dtw_loss_functions.otw}}\label{\detokenize{dtw_loss_functions.otw:otw}}\label{\detokenize{dtw_loss_functions.otw::doc}}\index{module@\spxentry{module}!dtw\_loss\_functions.otw@\spxentry{dtw\_loss\_functions.otw}}\index{dtw\_loss\_functions.otw@\spxentry{dtw\_loss\_functions.otw}!module@\spxentry{module}}
\sphinxAtStartPar
Implementation of OTW, from “OTW: Optimal Transport Warping for Time Series” (by Fabian Latorre, Chenghao Liu, Doyen Sahoo, and Steven C.H. Hoi)

\sphinxAtStartPar
For more details see :
\sphinxhyphen{} \sphinxurl{https://ieeexplore.ieee.org/document/10095915}
\sphinxhyphen{} \sphinxurl{https://arxiv.org/abs/2306.00620}


\section{Authors}
\label{\detokenize{dtw_loss_functions.otw:authors}}
\sphinxAtStartPar
Alberto Zancanaro \textless{}\sphinxhref{mailto:alberto.zancanaro@uni.lu}{alberto.zancanaro@uni.lu}\textgreater{}
\index{otw (class in dtw\_loss\_functions.otw)@\spxentry{otw}\spxextra{class in dtw\_loss\_functions.otw}}

\begin{fulllineitems}
\phantomsection\label{\detokenize{dtw_loss_functions.otw:dtw_loss_functions.otw.otw}}
\pysigstartsignatures
\pysiglinewithargsret
{\sphinxbfcode{\sphinxupquote{\DUrole{k}{class}\DUrole{w}{ }}}\sphinxcode{\sphinxupquote{dtw\_loss\_functions.otw.}}\sphinxbfcode{\sphinxupquote{otw}}}
{\sphinxparam{\DUrole{n}{m}\DUrole{p}{:}\DUrole{w}{ }\DUrole{n}{float}\DUrole{w}{ }\DUrole{o}{=}\DUrole{w}{ }\DUrole{default_value}{1}}\sphinxparamcomma \sphinxparam{\DUrole{n}{s}\DUrole{p}{:}\DUrole{w}{ }\DUrole{n}{int\DUrole{w}{ }\DUrole{p}{|}\DUrole{w}{ }float}\DUrole{w}{ }\DUrole{o}{=}\DUrole{w}{ }\DUrole{default_value}{0.5}}\sphinxparamcomma \sphinxparam{\DUrole{n}{beta}\DUrole{p}{:}\DUrole{w}{ }\DUrole{n}{float}\DUrole{w}{ }\DUrole{o}{=}\DUrole{w}{ }\DUrole{default_value}{1}}\sphinxparamcomma \sphinxparam{\DUrole{n}{reduction}\DUrole{p}{:}\DUrole{w}{ }\DUrole{n}{str}\DUrole{w}{ }\DUrole{o}{=}\DUrole{w}{ }\DUrole{default_value}{\textquotesingle{}mean\textquotesingle{}}}}
{}
\pysigstopsignatures
\sphinxAtStartPar
Bases: \sphinxcode{\sphinxupquote{object}}
\subsubsection*{Methods}


\begin{savenotes}
\sphinxatlongtablestart
\sphinxthistablewithglobalstyle
\sphinxthistablewithnovlinesstyle
\makeatletter
  \LTleft \@totalleftmargin plus1fill
  \LTright\dimexpr\columnwidth-\@totalleftmargin-\linewidth\relax plus1fill
\makeatother
\begin{longtable}{\X{1}{2}\X{1}{2}}
\sphinxtoprule
\endfirsthead

\multicolumn{2}{c}{\sphinxnorowcolor
    \makebox[0pt]{\sphinxtablecontinued{\tablename\ \thetable{} \textendash{} continued from previous page}}%
}\\
\sphinxtoprule
\endhead

\sphinxbottomrule
\multicolumn{2}{r}{\sphinxnorowcolor
    \makebox[0pt][r]{\sphinxtablecontinued{continues on next page}}%
}\\
\endfoot

\endlastfoot
\sphinxtableatstartofbodyhook
\begin{varwidth}[t]{\sphinxcolwidth{1}{2}}
\sphinxAtStartPar
\sphinxcode{\sphinxupquote{\_\_call\_\_}}(x, y)
\sphinxbeforeendvarwidth
\end{varwidth}%
&\begin{varwidth}[t]{\sphinxcolwidth{1}{2}}
\sphinxAtStartPar
Computes the OTW distance between two time series.
\sphinxbeforeendvarwidth
\end{varwidth}%
\\
\sphinxbottomrule
\end{longtable}
\sphinxtableafterendhook
\sphinxatlongtableend
\end{savenotes}

\end{fulllineitems}

\index{otw\_distance() (in module dtw\_loss\_functions.otw)@\spxentry{otw\_distance()}\spxextra{in module dtw\_loss\_functions.otw}}

\begin{fulllineitems}
\phantomsection\label{\detokenize{dtw_loss_functions.otw:dtw_loss_functions.otw.otw_distance}}
\pysigstartsignatures
\pysiglinewithargsret
{\sphinxcode{\sphinxupquote{dtw\_loss\_functions.otw.}}\sphinxbfcode{\sphinxupquote{otw\_distance}}}
{\sphinxparam{\DUrole{n}{x}\DUrole{p}{:}\DUrole{w}{ }\DUrole{n}{Tensor}}\sphinxparamcomma \sphinxparam{\DUrole{n}{y}\DUrole{p}{:}\DUrole{w}{ }\DUrole{n}{Tensor}}\sphinxparamcomma \sphinxparam{\DUrole{n}{m}\DUrole{p}{:}\DUrole{w}{ }\DUrole{n}{float}\DUrole{w}{ }\DUrole{o}{=}\DUrole{w}{ }\DUrole{default_value}{1}}\sphinxparamcomma \sphinxparam{\DUrole{n}{s}\DUrole{p}{:}\DUrole{w}{ }\DUrole{n}{int\DUrole{w}{ }\DUrole{p}{|}\DUrole{w}{ }float}\DUrole{w}{ }\DUrole{o}{=}\DUrole{w}{ }\DUrole{default_value}{0.5}}\sphinxparamcomma \sphinxparam{\DUrole{n}{beta}\DUrole{p}{:}\DUrole{w}{ }\DUrole{n}{float}\DUrole{w}{ }\DUrole{o}{=}\DUrole{w}{ }\DUrole{default_value}{1}}\sphinxparamcomma \sphinxparam{\DUrole{n}{reduction}\DUrole{p}{:}\DUrole{w}{ }\DUrole{n}{str}\DUrole{w}{ }\DUrole{o}{=}\DUrole{w}{ }\DUrole{default_value}{\textquotesingle{}mean\textquotesingle{}}}}
{{ $\rightarrow$ Tensor}}
\pysigstopsignatures
\sphinxAtStartPar
Implements the OTW distance between two time series, as defined in equation (10) of the paper.
\begin{quote}\begin{description}
\sphinxlineitem{Parameters}\begin{itemize}
\item {} 
\sphinxAtStartPar
\sphinxstyleliteralstrong{\sphinxupquote{x}} (\sphinxstyleliteralemphasis{\sphinxupquote{torch.Tensor}}) \textendash{} First time series, of shape (B, L) where B is the batch size and L is the length of the time series.

\item {} 
\sphinxAtStartPar
\sphinxstyleliteralstrong{\sphinxupquote{y}} (\sphinxstyleliteralemphasis{\sphinxupquote{torch.Tensor}}) \textendash{} Second time series, of shape (B, L) where B is the batch size and L is the length of the time series.

\item {} 
\sphinxAtStartPar
\sphinxstyleliteralstrong{\sphinxupquote{m}} (\sphinxstyleliteralemphasis{\sphinxupquote{float}}) \textendash{} Waste cost parameter, default is 1.

\item {} 
\sphinxAtStartPar
\sphinxstyleliteralstrong{\sphinxupquote{s}} (\sphinxstyleliteralemphasis{\sphinxupquote{int}}\sphinxstyleliteralemphasis{\sphinxupquote{ | }}\sphinxstyleliteralemphasis{\sphinxupquote{float}}) \textendash{} Window size parameter, it can be an integer or a float between 0 and 1. Default is 0.5.
If float, it is interpreted as a fraction of the length of the time series.
If integer, it is interpreted as the number of time steps.

\item {} 
\sphinxAtStartPar
\sphinxstyleliteralstrong{\sphinxupquote{beta}} (\sphinxstyleliteralemphasis{\sphinxupquote{float}}) \textendash{} Hyperparameter for the smooth l1 loss, default is 1.

\end{itemize}

\sphinxlineitem{Returns}
\sphinxAtStartPar
OTW distance between the two time series

\sphinxlineitem{Return type}
\sphinxAtStartPar
torch.Tensor

\end{description}\end{quote}

\end{fulllineitems}

\index{smooth\_l1\_loss() (in module dtw\_loss\_functions.otw)@\spxentry{smooth\_l1\_loss()}\spxextra{in module dtw\_loss\_functions.otw}}

\begin{fulllineitems}
\phantomsection\label{\detokenize{dtw_loss_functions.otw:dtw_loss_functions.otw.smooth_l1_loss}}
\pysigstartsignatures
\pysiglinewithargsret
{\sphinxcode{\sphinxupquote{dtw\_loss\_functions.otw.}}\sphinxbfcode{\sphinxupquote{smooth\_l1\_loss}}}
{\sphinxparam{\DUrole{n}{x}\DUrole{p}{:}\DUrole{w}{ }\DUrole{n}{Tensor}}\sphinxparamcomma \sphinxparam{\DUrole{n}{beta}\DUrole{p}{:}\DUrole{w}{ }\DUrole{n}{float}}\sphinxparamcomma \sphinxparam{\DUrole{n}{reduction}\DUrole{o}{=}\DUrole{default_value}{\textquotesingle{}mean\textquotesingle{}}}}
{{ $\rightarrow$ Tensor}}
\pysigstopsignatures
\sphinxAtStartPar
Computes the smooth l1 of the input tensor x, as defined in equation (9) of the paper.
\begin{quote}\begin{description}
\sphinxlineitem{Parameters}\begin{itemize}
\item {} 
\sphinxAtStartPar
\sphinxstyleliteralstrong{\sphinxupquote{x}} (\sphinxstyleliteralemphasis{\sphinxupquote{torch.Tensor}}) \textendash{} Input tensor of shape (B). Each element corresponds to the difference between two time series.

\item {} 
\sphinxAtStartPar
\sphinxstyleliteralstrong{\sphinxupquote{beta}} (\sphinxstyleliteralemphasis{\sphinxupquote{float}}) \textendash{} Hyperparameter for the smooth l1 loss.

\item {} 
\sphinxAtStartPar
\sphinxstyleliteralstrong{\sphinxupquote{reduction}} (\sphinxstyleliteralemphasis{\sphinxupquote{str}}) \textendash{} Specifies the reduction to apply to the output: ‘none’ | ‘mean’ | ‘sum’. Default: ‘mean’.

\end{itemize}

\sphinxlineitem{Returns}
\sphinxAtStartPar
Smooth l1 loss of the input tensor. If reduction is ‘none’, the output has the same shape as x. If reduction is ‘mean’ or ‘sum’, the output is a scalar.

\sphinxlineitem{Return type}
\sphinxAtStartPar
torch.Tensor

\end{description}\end{quote}

\end{fulllineitems}

\index{window\_cumsum() (in module dtw\_loss\_functions.otw)@\spxentry{window\_cumsum()}\spxextra{in module dtw\_loss\_functions.otw}}

\begin{fulllineitems}
\phantomsection\label{\detokenize{dtw_loss_functions.otw:dtw_loss_functions.otw.window_cumsum}}
\pysigstartsignatures
\pysiglinewithargsret
{\sphinxcode{\sphinxupquote{dtw\_loss\_functions.otw.}}\sphinxbfcode{\sphinxupquote{window\_cumsum}}}
{\sphinxparam{\DUrole{n}{x}\DUrole{p}{:}\DUrole{w}{ }\DUrole{n}{Tensor}}\sphinxparamcomma \sphinxparam{\DUrole{n}{s}\DUrole{p}{:}\DUrole{w}{ }\DUrole{n}{int}}}
{{ $\rightarrow$ Tensor}}
\pysigstopsignatures
\sphinxAtStartPar
Computes the cumulative sum of the input tensor x as defined in equation (7) of the paper.

\sphinxAtStartPar
Given a time series A represented as an array of values \sphinxcode{\sphinxupquote{{[}a1, a2, ..., aL{]}}}, the window cumsum is computed as :
\sphinxcode{\sphinxupquote{window\_cumsum(A) = cumsum(A) \sphinxhyphen{} cumsum(A{[}0:L\sphinxhyphen{}s{]})}}
(i.e. the cumsum of all the array minus the cumsum of the array excluding the last s elements)
\begin{quote}\begin{description}
\sphinxlineitem{Parameters}\begin{itemize}
\item {} 
\sphinxAtStartPar
\sphinxstyleliteralstrong{\sphinxupquote{x}} (\sphinxstyleliteralemphasis{\sphinxupquote{torch.Tensor}}) \textendash{} Input tensor of shape (B, L).

\item {} 
\sphinxAtStartPar
\sphinxstyleliteralstrong{\sphinxupquote{s}} (\sphinxstyleliteralemphasis{\sphinxupquote{int}}) \textendash{} Window size.

\end{itemize}

\sphinxlineitem{Returns}
\sphinxAtStartPar
Cumulative sum over the sliding window, of shape (B).

\sphinxlineitem{Return type}
\sphinxAtStartPar
torch.Tensor

\end{description}\end{quote}

\end{fulllineitems}


\sphinxstepscope


\chapter{Soft DTW}
\label{\detokenize{dtw_loss_functions.soft_dtw_cuda:module-dtw_loss_functions.soft_dtw_cuda}}\label{\detokenize{dtw_loss_functions.soft_dtw_cuda:soft-dtw}}\label{\detokenize{dtw_loss_functions.soft_dtw_cuda::doc}}\index{module@\spxentry{module}!dtw\_loss\_functions.soft\_dtw\_cuda@\spxentry{dtw\_loss\_functions.soft\_dtw\_cuda}}\index{dtw\_loss\_functions.soft\_dtw\_cuda@\spxentry{dtw\_loss\_functions.soft\_dtw\_cuda}!module@\spxentry{module}}
\sphinxAtStartPar
Cuda Implementation of SoftDTW by Mehran Maghoumi.

\sphinxAtStartPar
The code here is identical to the one in the original repository (\sphinxurl{https://github.com/Maghoumi/pytorch-softdtw-cuda}).
The only changes are the modification of the docstring to numpydoc style, to make it consistent with the rest of the codebase.
\subsubsection*{Example}

\begin{sphinxVerbatim}[commandchars=\\\{\}]
\PYG{g+gp}{\PYGZgt{}\PYGZgt{}\PYGZgt{} }\PYG{k+kn}{from}\PYG{+w}{ }\PYG{n+nn}{dtw\PYGZus{}loss\PYGZus{}functions}\PYG{+w}{ }\PYG{k+kn}{import} \PYG{n}{soft\PYGZus{}dtw\PYGZus{}cuda}
\PYG{g+gp}{\PYGZgt{}\PYGZgt{}\PYGZgt{} }\PYG{k+kn}{import}\PYG{+w}{ }\PYG{n+nn}{torch}
\PYG{g+gp}{\PYGZgt{}\PYGZgt{}\PYGZgt{} }\PYG{n}{use\PYGZus{}cuda} \PYG{o}{=} \PYG{n}{torch}\PYG{o}{.}\PYG{n}{cuda}\PYG{o}{.}\PYG{n}{is\PYGZus{}available}\PYG{p}{(}\PYG{p}{)}
\PYG{g+gp}{\PYGZgt{}\PYGZgt{}\PYGZgt{} }\PYG{n}{device} \PYG{o}{=} \PYG{l+s+s1}{\PYGZsq{}}\PYG{l+s+s1}{cuda}\PYG{l+s+s1}{\PYGZsq{}} \PYG{k}{if} \PYG{n}{use\PYGZus{}cuda} \PYG{k}{else} \PYG{l+s+s1}{\PYGZsq{}}\PYG{l+s+s1}{cpu}\PYG{l+s+s1}{\PYGZsq{}}
\PYG{g+gp}{\PYGZgt{}\PYGZgt{}\PYGZgt{} }\PYG{n}{batch\PYGZus{}size} \PYG{o}{=} \PYG{l+m+mi}{5}
\PYG{g+gp}{\PYGZgt{}\PYGZgt{}\PYGZgt{} }\PYG{n}{time\PYGZus{}samples} \PYG{o}{=} \PYG{l+m+mi}{300}
\PYG{g+gp}{\PYGZgt{}\PYGZgt{}\PYGZgt{} }\PYG{n}{channels} \PYG{o}{=} \PYG{l+m+mi}{1}
\PYG{g+gp}{\PYGZgt{}\PYGZgt{}\PYGZgt{} }\PYG{n}{x}   \PYG{o}{=} \PYG{n}{torch}\PYG{o}{.}\PYG{n}{randn}\PYG{p}{(}\PYG{n}{batch\PYGZus{}size}\PYG{p}{,} \PYG{n}{time\PYGZus{}samples}\PYG{p}{,} \PYG{n}{channels}\PYG{p}{)}\PYG{o}{.}\PYG{n}{to}\PYG{p}{(}\PYG{n}{device}\PYG{p}{)}
\PYG{g+gp}{\PYGZgt{}\PYGZgt{}\PYGZgt{} }\PYG{n}{x\PYGZus{}r} \PYG{o}{=} \PYG{n}{torch}\PYG{o}{.}\PYG{n}{randn}\PYG{p}{(}\PYG{n}{batch\PYGZus{}size}\PYG{p}{,} \PYG{n}{time\PYGZus{}samples}\PYG{p}{,} \PYG{n}{channels}\PYG{p}{)}\PYG{o}{.}\PYG{n}{to}\PYG{p}{(}\PYG{n}{device}\PYG{p}{)}
\PYG{g+gp}{\PYGZgt{}\PYGZgt{}\PYGZgt{} }\PYG{n}{sdtw\PYGZus{}loss} \PYG{o}{=} \PYG{n}{soft\PYGZus{}dtw\PYGZus{}cuda}\PYG{o}{.}\PYG{n}{SoftDTW}\PYG{p}{(}\PYG{n}{use\PYGZus{}cuda} \PYG{o}{=} \PYG{n}{use\PYGZus{}cuda}\PYG{p}{)}
\PYG{g+gp}{\PYGZgt{}\PYGZgt{}\PYGZgt{} }\PYG{n}{output\PYGZus{}sdtw} \PYG{o}{=} \PYG{n}{sdtw\PYGZus{}loss}\PYG{p}{(}\PYG{n}{x}\PYG{p}{,} \PYG{n}{x\PYGZus{}r}\PYG{p}{)}
\end{sphinxVerbatim}


\section{Authors}
\label{\detokenize{dtw_loss_functions.soft_dtw_cuda:authors}}
\sphinxAtStartPar
Mehran Maghoumi
\index{SoftDTW (class in dtw\_loss\_functions.soft\_dtw\_cuda)@\spxentry{SoftDTW}\spxextra{class in dtw\_loss\_functions.soft\_dtw\_cuda}}

\begin{fulllineitems}
\phantomsection\label{\detokenize{dtw_loss_functions.soft_dtw_cuda:dtw_loss_functions.soft_dtw_cuda.SoftDTW}}
\pysigstartsignatures
\pysiglinewithargsret
{\sphinxbfcode{\sphinxupquote{\DUrole{k}{class}\DUrole{w}{ }}}\sphinxcode{\sphinxupquote{dtw\_loss\_functions.soft\_dtw\_cuda.}}\sphinxbfcode{\sphinxupquote{SoftDTW}}}
{\sphinxparam{\DUrole{n}{use\_cuda}}\sphinxparamcomma \sphinxparam{\DUrole{n}{gamma}\DUrole{o}{=}\DUrole{default_value}{1.0}}\sphinxparamcomma \sphinxparam{\DUrole{n}{normalize}\DUrole{o}{=}\DUrole{default_value}{False}}\sphinxparamcomma \sphinxparam{\DUrole{n}{bandwidth}\DUrole{o}{=}\DUrole{default_value}{None}}\sphinxparamcomma \sphinxparam{\DUrole{n}{dist\_func}\DUrole{o}{=}\DUrole{default_value}{None}}}
{}
\pysigstopsignatures
\sphinxAtStartPar
Bases: \sphinxcode{\sphinxupquote{Module}}

\sphinxAtStartPar
The soft DTW implementation that optionally supports CUDA
\index{normalize (dtw\_loss\_functions.soft\_dtw\_cuda.SoftDTW attribute)@\spxentry{normalize}\spxextra{dtw\_loss\_functions.soft\_dtw\_cuda.SoftDTW attribute}}

\begin{fulllineitems}
\phantomsection\label{\detokenize{dtw_loss_functions.soft_dtw_cuda:dtw_loss_functions.soft_dtw_cuda.SoftDTW.normalize}}
\pysigstartsignatures
\pysigline
{\sphinxbfcode{\sphinxupquote{normalize}}}
\pysigstopsignatures
\sphinxAtStartPar
Flag indicating whether to perform normalization (as discussed in discussed in \sphinxurl{https://github.com/mblondel/soft-dtw/issues/10\#issuecomment-383564790})
Note that if normalize is set to True, the SoftDTW divergence will be computed, which is defined as \sphinxcode{\sphinxupquote{SoftDTW(X, Y) \sphinxhyphen{} 1/2 * (SoftDTW(X, X) + SoftDTW(Y, Y))}}.
\begin{quote}\begin{description}
\sphinxlineitem{Type}
\sphinxAtStartPar
bool

\end{description}\end{quote}

\end{fulllineitems}

\index{gamma (dtw\_loss\_functions.soft\_dtw\_cuda.SoftDTW attribute)@\spxentry{gamma}\spxextra{dtw\_loss\_functions.soft\_dtw\_cuda.SoftDTW attribute}}

\begin{fulllineitems}
\phantomsection\label{\detokenize{dtw_loss_functions.soft_dtw_cuda:dtw_loss_functions.soft_dtw_cuda.SoftDTW.gamma}}
\pysigstartsignatures
\pysigline
{\sphinxbfcode{\sphinxupquote{gamma}}}
\pysigstopsignatures
\sphinxAtStartPar
SoftDTW’s gamma parameter
\begin{quote}\begin{description}
\sphinxlineitem{Type}
\sphinxAtStartPar
float

\end{description}\end{quote}

\end{fulllineitems}

\index{bandwidth (dtw\_loss\_functions.soft\_dtw\_cuda.SoftDTW attribute)@\spxentry{bandwidth}\spxextra{dtw\_loss\_functions.soft\_dtw\_cuda.SoftDTW attribute}}

\begin{fulllineitems}
\phantomsection\label{\detokenize{dtw_loss_functions.soft_dtw_cuda:dtw_loss_functions.soft_dtw_cuda.SoftDTW.bandwidth}}
\pysigstartsignatures
\pysigline
{\sphinxbfcode{\sphinxupquote{bandwidth}}}
\pysigstopsignatures
\sphinxAtStartPar
Sakoe\sphinxhyphen{}Chiba bandwidth for pruning. Passing ‘None’ will disable pruning.
\begin{quote}\begin{description}
\sphinxlineitem{Type}
\sphinxAtStartPar
float or None

\end{description}\end{quote}

\end{fulllineitems}

\index{dist\_func (dtw\_loss\_functions.soft\_dtw\_cuda.SoftDTW attribute)@\spxentry{dist\_func}\spxextra{dtw\_loss\_functions.soft\_dtw\_cuda.SoftDTW attribute}}

\begin{fulllineitems}
\phantomsection\label{\detokenize{dtw_loss_functions.soft_dtw_cuda:dtw_loss_functions.soft_dtw_cuda.SoftDTW.dist_func}}
\pysigstartsignatures
\pysigline
{\sphinxbfcode{\sphinxupquote{dist\_func}}}
\pysigstopsignatures
\sphinxAtStartPar
Optional point\sphinxhyphen{}wise distance function to use. If ‘None’, then a default Euclidean distance function will be used.
\begin{quote}\begin{description}
\sphinxlineitem{Type}
\sphinxAtStartPar
function or None

\end{description}\end{quote}

\end{fulllineitems}

\begin{quote}\begin{description}
\sphinxlineitem{Parameters}\begin{itemize}
\item {} 
\sphinxAtStartPar
\sphinxstyleliteralstrong{\sphinxupquote{use\_cuda}} (\sphinxstyleliteralemphasis{\sphinxupquote{bool}}) \textendash{} Flag indicating whether the CUDA implementation should be used

\item {} 
\sphinxAtStartPar
\sphinxstyleliteralstrong{\sphinxupquote{gamma}} (\sphinxstyleliteralemphasis{\sphinxupquote{float}}) \textendash{} SoftDTW’s gamma parameter

\item {} 
\sphinxAtStartPar
\sphinxstyleliteralstrong{\sphinxupquote{normalize}} (\sphinxstyleliteralemphasis{\sphinxupquote{bool}}) \textendash{} Flag indicating whether to perform normalization (as discussed in discussed in \sphinxurl{https://github.com/mblondel/soft-dtw/issues/10\#issuecomment-383564790})
Note that if normalize is set to True, the SoftDTW divergence will be computed, which is defined as \sphinxcode{\sphinxupquote{SoftDTW(X, Y) \sphinxhyphen{} 1/2 * (SoftDTW(X, X) + SoftDTW(Y, Y))}}.

\item {} 
\sphinxAtStartPar
\sphinxstyleliteralstrong{\sphinxupquote{bandwidth}} (\sphinxstyleliteralemphasis{\sphinxupquote{float}}\sphinxstyleliteralemphasis{\sphinxupquote{ or }}\sphinxstyleliteralemphasis{\sphinxupquote{None}}) \textendash{} Sakoe\sphinxhyphen{}Chiba bandwidth for pruning. Passing ‘None’ will disable pruning.

\item {} 
\sphinxAtStartPar
\sphinxstyleliteralstrong{\sphinxupquote{dist\_func}} (\sphinxstyleliteralemphasis{\sphinxupquote{function}}\sphinxstyleliteralemphasis{\sphinxupquote{ or }}\sphinxstyleliteralemphasis{\sphinxupquote{None}}) \textendash{} Optional point\sphinxhyphen{}wise distance function to use. If ‘None’, then a default Euclidean distance function will be used.

\end{itemize}

\end{description}\end{quote}
\subsubsection*{Methods}


\begin{savenotes}
\sphinxatlongtablestart
\sphinxthistablewithglobalstyle
\sphinxthistablewithnovlinesstyle
\makeatletter
  \LTleft \@totalleftmargin plus1fill
  \LTright\dimexpr\columnwidth-\@totalleftmargin-\linewidth\relax plus1fill
\makeatother
\begin{longtable}{\X{1}{2}\X{1}{2}}
\sphinxtoprule
\endfirsthead

\multicolumn{2}{c}{\sphinxnorowcolor
    \makebox[0pt]{\sphinxtablecontinued{\tablename\ \thetable{} \textendash{} continued from previous page}}%
}\\
\sphinxtoprule
\endhead

\sphinxbottomrule
\multicolumn{2}{r}{\sphinxnorowcolor
    \makebox[0pt][r]{\sphinxtablecontinued{continues on next page}}%
}\\
\endfoot

\endlastfoot
\sphinxtableatstartofbodyhook
\begin{varwidth}[t]{\sphinxcolwidth{1}{2}}
\sphinxAtStartPar
\sphinxcode{\sphinxupquote{forward}}(X, Y)
\sphinxbeforeendvarwidth
\end{varwidth}%
&\begin{varwidth}[t]{\sphinxcolwidth{1}{2}}
\sphinxAtStartPar
Compute the soft\sphinxhyphen{}DTW value between X and Y :param X: One batch of examples, batch\_size x seq\_len x dims :param Y: The other batch of examples, batch\_size x seq\_len x dims :return: The computed results
\sphinxbeforeendvarwidth
\end{varwidth}%
\\
\sphinxbottomrule
\end{longtable}
\sphinxtableafterendhook
\sphinxatlongtableend
\end{savenotes}
\index{forward() (dtw\_loss\_functions.soft\_dtw\_cuda.SoftDTW method)@\spxentry{forward()}\spxextra{dtw\_loss\_functions.soft\_dtw\_cuda.SoftDTW method}}

\begin{fulllineitems}
\phantomsection\label{\detokenize{dtw_loss_functions.soft_dtw_cuda:dtw_loss_functions.soft_dtw_cuda.SoftDTW.forward}}
\pysigstartsignatures
\pysiglinewithargsret
{\sphinxbfcode{\sphinxupquote{forward}}}
{\sphinxparam{\DUrole{n}{X}}\sphinxparamcomma \sphinxparam{\DUrole{n}{Y}}}
{}
\pysigstopsignatures
\sphinxAtStartPar
Compute the soft\sphinxhyphen{}DTW value between X and Y
:param X: One batch of examples, batch\_size x seq\_len x dims
:param Y: The other batch of examples, batch\_size x seq\_len x dims
:return: The computed results

\end{fulllineitems}


\end{fulllineitems}

\index{compute\_softdtw() (in module dtw\_loss\_functions.soft\_dtw\_cuda)@\spxentry{compute\_softdtw()}\spxextra{in module dtw\_loss\_functions.soft\_dtw\_cuda}}

\begin{fulllineitems}
\phantomsection\label{\detokenize{dtw_loss_functions.soft_dtw_cuda:dtw_loss_functions.soft_dtw_cuda.compute_softdtw}}
\pysigstartsignatures
\pysiglinewithargsret
{\sphinxcode{\sphinxupquote{dtw\_loss\_functions.soft\_dtw\_cuda.}}\sphinxbfcode{\sphinxupquote{compute\_softdtw}}}
{\sphinxparam{\DUrole{n}{D}}\sphinxparamcomma \sphinxparam{\DUrole{n}{gamma}}\sphinxparamcomma \sphinxparam{\DUrole{n}{bandwidth}}}
{}
\pysigstopsignatures
\end{fulllineitems}

\index{compute\_softdtw\_backward() (in module dtw\_loss\_functions.soft\_dtw\_cuda)@\spxentry{compute\_softdtw\_backward()}\spxextra{in module dtw\_loss\_functions.soft\_dtw\_cuda}}

\begin{fulllineitems}
\phantomsection\label{\detokenize{dtw_loss_functions.soft_dtw_cuda:dtw_loss_functions.soft_dtw_cuda.compute_softdtw_backward}}
\pysigstartsignatures
\pysiglinewithargsret
{\sphinxcode{\sphinxupquote{dtw\_loss\_functions.soft\_dtw\_cuda.}}\sphinxbfcode{\sphinxupquote{compute\_softdtw\_backward}}}
{\sphinxparam{\DUrole{n}{D\_}}\sphinxparamcomma \sphinxparam{\DUrole{n}{R}}\sphinxparamcomma \sphinxparam{\DUrole{n}{gamma}}\sphinxparamcomma \sphinxparam{\DUrole{n}{bandwidth}}}
{}
\pysigstopsignatures
\end{fulllineitems}

\index{compute\_softdtw\_backward\_cuda() (in module dtw\_loss\_functions.soft\_dtw\_cuda)@\spxentry{compute\_softdtw\_backward\_cuda()}\spxextra{in module dtw\_loss\_functions.soft\_dtw\_cuda}}

\begin{fulllineitems}
\phantomsection\label{\detokenize{dtw_loss_functions.soft_dtw_cuda:dtw_loss_functions.soft_dtw_cuda.compute_softdtw_backward_cuda}}
\pysigstartsignatures
\pysiglinewithargsret
{\sphinxcode{\sphinxupquote{dtw\_loss\_functions.soft\_dtw\_cuda.}}\sphinxbfcode{\sphinxupquote{compute\_softdtw\_backward\_cuda}}}
{\sphinxparam{\DUrole{n}{D}}\sphinxparamcomma \sphinxparam{\DUrole{n}{R}}\sphinxparamcomma \sphinxparam{\DUrole{n}{inv\_gamma}}\sphinxparamcomma \sphinxparam{\DUrole{n}{bandwidth}}\sphinxparamcomma \sphinxparam{\DUrole{n}{max\_i}}\sphinxparamcomma \sphinxparam{\DUrole{n}{max\_j}}\sphinxparamcomma \sphinxparam{\DUrole{n}{n\_passes}}\sphinxparamcomma \sphinxparam{\DUrole{n}{E}}}
{}
\pysigstopsignatures
\end{fulllineitems}

\index{compute\_softdtw\_cuda() (in module dtw\_loss\_functions.soft\_dtw\_cuda)@\spxentry{compute\_softdtw\_cuda()}\spxextra{in module dtw\_loss\_functions.soft\_dtw\_cuda}}

\begin{fulllineitems}
\phantomsection\label{\detokenize{dtw_loss_functions.soft_dtw_cuda:dtw_loss_functions.soft_dtw_cuda.compute_softdtw_cuda}}
\pysigstartsignatures
\pysiglinewithargsret
{\sphinxcode{\sphinxupquote{dtw\_loss\_functions.soft\_dtw\_cuda.}}\sphinxbfcode{\sphinxupquote{compute\_softdtw\_cuda}}}
{\sphinxparam{\DUrole{n}{D}}\sphinxparamcomma \sphinxparam{\DUrole{n}{gamma}}\sphinxparamcomma \sphinxparam{\DUrole{n}{bandwidth}}\sphinxparamcomma \sphinxparam{\DUrole{n}{max\_i}}\sphinxparamcomma \sphinxparam{\DUrole{n}{max\_j}}\sphinxparamcomma \sphinxparam{\DUrole{n}{n\_passes}}\sphinxparamcomma \sphinxparam{\DUrole{n}{R}}}
{}
\pysigstopsignatures
\sphinxAtStartPar
Computes the soft\sphinxhyphen{}DTW value between two sequences using CUDA.

\end{fulllineitems}

\index{profile() (in module dtw\_loss\_functions.soft\_dtw\_cuda)@\spxentry{profile()}\spxextra{in module dtw\_loss\_functions.soft\_dtw\_cuda}}

\begin{fulllineitems}
\phantomsection\label{\detokenize{dtw_loss_functions.soft_dtw_cuda:dtw_loss_functions.soft_dtw_cuda.profile}}
\pysigstartsignatures
\pysiglinewithargsret
{\sphinxcode{\sphinxupquote{dtw\_loss\_functions.soft\_dtw\_cuda.}}\sphinxbfcode{\sphinxupquote{profile}}}
{\sphinxparam{\DUrole{n}{batch\_size}}\sphinxparamcomma \sphinxparam{\DUrole{n}{seq\_len\_a}}\sphinxparamcomma \sphinxparam{\DUrole{n}{seq\_len\_b}}\sphinxparamcomma \sphinxparam{\DUrole{n}{dims}}\sphinxparamcomma \sphinxparam{\DUrole{n}{tol\_backward}}}
{}
\pysigstopsignatures
\end{fulllineitems}

\index{timed\_run() (in module dtw\_loss\_functions.soft\_dtw\_cuda)@\spxentry{timed\_run()}\spxextra{in module dtw\_loss\_functions.soft\_dtw\_cuda}}

\begin{fulllineitems}
\phantomsection\label{\detokenize{dtw_loss_functions.soft_dtw_cuda:dtw_loss_functions.soft_dtw_cuda.timed_run}}
\pysigstartsignatures
\pysiglinewithargsret
{\sphinxcode{\sphinxupquote{dtw\_loss\_functions.soft\_dtw\_cuda.}}\sphinxbfcode{\sphinxupquote{timed\_run}}}
{\sphinxparam{\DUrole{n}{a}}\sphinxparamcomma \sphinxparam{\DUrole{n}{b}}\sphinxparamcomma \sphinxparam{\DUrole{n}{sdtw}}}
{}
\pysigstopsignatures
\sphinxAtStartPar
Runs a and b through sdtw, and times the forward and backward passes.
Assumes that a requires gradients.
:return: timing, forward result, backward result

\end{fulllineitems}



\chapter{Indices and Tables}
\label{\detokenize{index:indices-and-tables}}\begin{itemize}
\item {} 
\sphinxAtStartPar
\DUrole{xref}{\DUrole{std}{\DUrole{std-ref}{genindex}}}

\item {} 
\sphinxAtStartPar
\DUrole{xref}{\DUrole{std}{\DUrole{std-ref}{modindex}}}

\item {} 
\sphinxAtStartPar
\DUrole{xref}{\DUrole{std}{\DUrole{std-ref}{search}}}

\end{itemize}



\renewcommand{\indexname}{Index}
\printindex
\end{document}